\section{Introduction}

An issued wind power forecast defines a production target, but the reserve request must still account for possible shortfall. Measurements arriving between forecast issuance and the reserve decision can update that request. The practical question is which measurements provide enough additional value to justify their inclusion. This question concerns the information supplied to the reserve layer after the forecast has already been produced.

Previous studies establish several links between wind observations, forecasts and operating decisions. Conditional forecast-error distributions support probabilistic reserve sizing \cite{bruninx2014}, and turbine-level data improve probabilistic wind farm forecasts \cite{gilbert2020}. Flexible density models further describe uncertainty across forecast horizons \cite{wen2022}. Offering strategies and capacity-firming methods evaluate forecasts through reserve penalties or operating returns \cite{soares2016,dumas2022}. These contributions motivate assessing information by its consequences for the decision that uses it.

Input selection based on downstream cost also has established precedents. Feature-driven renewable trading and newsvendor models select or transform predictors according to decision loss \cite{munoz2020featuredriven,serrano2024featureselection}. Missing-feature formulations consider inputs that become unavailable during operation \cite{kuhnau2023missing}. The remaining question here is more specific. Given an unchanged forecast, how much telemetry must be supplied separately to the reserve learner? The answer depends on what the forecast already summarizes and what the additional measurement contributes at decision time.

Current active power and blade pitch provide a useful matched comparison. Active power measures electrical output, whereas pitch records blade actuation. Short-term quantile models use recent production to characterize wind uncertainty \cite{pritchard2011shortterm}. This motivates testing the incremental decision value of current power against a contemporaneous alternative under the same forecast and telemetry background. Matching the added input width separates this comparison from the benefit of supplying more derived features.

Measurement value alone does not establish how far the input set can be reduced. A compact rule must also retain acceptable cost and shortfall frequency when applied to later periods. We address these two questions using a common forecast commitment and an asymmetric shortfall loss. The first study compares single-column power and pitch updates on data from 134 turbines. The second tests a historically selected current-power rule against a three-channel history reference at Hill of Towie in 2025 and January--April 2026.

The contribution is a decision-based assessment of the telemetry interface between an existing forecast and a reserve request. The matched study quantifies the value of current power independently of two additional power-derived features. The subsequent-year study tests whether five raw reserve inputs retain the specified performance relative to 112 inputs. A secondary finite-family analysis asks the stronger question of whether all strictly smaller registered candidates can be excluded. Together, the studies connect measurement value to a concrete input-reduction decision.
